\loai{Vận dụng cơ bản định luật Charles}
\begin{vidu}%[1P5-4.2B][Loai1]
	Ở nhiệt độ 273\doC thể tích của một khối khí là 10 lít. Khi áp suất không đổi, thể tích của khí đó ở 546\doC là
	\choice
	{20 lít}
	{\True 15 lít}
	{12 lít}
	{13,5 lít}
	\loigiai{
		Áp dụng định luật Charles ta có: $\dfrac{V_2}{V_1}=\dfrac{T_2}{T_1}\Rightarrow V_2=\dfrac{V_1T_2}{T_1}=10.\dfrac{546+273}{273+273}=15\ \text{lít}$.
		\begin{note}
			Nhớ đổi đơn vị độ C ra độ K.
		\end{note}
	}
\end{vidu}
\loai{Thể tích hoặc nhiệt độ tăng (giảm) một lượng}
\begin{vidu}%[1P5-4.2K][Loai2]
	Đun nóng đẳng áp một khối khí lên đến $47\doC$ thì thể tích tăng thêm $\dfrac{1}{10}$ thể tích ban đầu. Nhiệt độ ban đầu là
	\choice
	{\True $17,9^\circ C$}
	{$117,9^\circ C$}
	{$19,7^\circ C$}
	{$47,9^\circ C$}
	\loigiai{
		\Binhluan{
			\begin{itemize}
				\item Áp dụng định luật Charles: $\dfrac{V_1}{T_1}=\dfrac{V_2}{T_2}\Rightarrow T_1=\dfrac{V_1T_2}{V_2}$.
				\item Ta có: $\left\{\begin{aligned}& V_2=V_1+\dfrac{V_1}{10}=1,1V_1 \\
					& T_2=273+47=320\ K 
				\end{aligned}\right.$
				\item $T_1 = \dfrac{320V_1}{1,1V_1}=\dfrac{320}{1,1}=290,9\ K \Rightarrow t_1 =290,9-273= 17,9\doC$.
			\end{itemize}
		}
		{Tính toán ở độ K rồi cuối cùng ta đổi ra độ C để kết luận}
	}
\end{vidu}
\loai{Bài toán liên quan khối lượng riêng}
\begin{vidu}%[1P5-4.2K][Loai3]
	Khối lượng riêng của một lượng chất khí trong bình A đậy kín ở $17\doC$ lớn hơn khối lượng riêng của cùng một lượng chất khí đó trong bình B đậy kín ở $37\doC$ bao nhiêu lần? Biết rằng áp suất của khí trong hai trường hợp bằng nhau.
	\choice
	{1,55}
	{7,11}
	{\True 1,07}
	{2,23}
	\loigiai{
		\Binhluan{
			\begin{itemize}
				\item Ở $T_1=290\ K$, khối lượng riêng của khí:  $D_1=\dfrac{m}{V_1}$.
				\item Ở $T_2=310\ K$, khối lượng riêng của khí:  $D_2=\dfrac{m}{V_2}$.
				\item Lập tỉ số ta có: $\dfrac{D_1}{D_2}=\dfrac{V_2}{V_1}=\dfrac{T_2}{T_1}=\dfrac{310}{290}=1,07$.
			\end{itemize}
		}
		{Mấu chốt là "cùng một lượng chất khí" nên $m=const$}
	}
\end{vidu}
\loai{Công chất khí thực hiện}
\begin{vidu}%[1P6-3.2K][Loai7]
	Một lượng khí lí tưởng chứa trong một xilanh có pit-tông chuyển động được. Các thông số trạng thái ban đầu của khí là $10\ dm^3$; 100 kPa; 300 K. Khí được làm lạnh theo một quá trình đẳng áp tới khi thể tích còn $6\ dm^3$. Xác định nhiệt độ cuối cùng của khí và tính công mà chất khí thực hiện được.
	\loigiai{
		\begin{itemize}
			\item Nhiệt độ cuối: $T_2=\dfrac{V_2T_1}{V_1}=180\ K$.
			\item Công chất khí thực hiện được: $A'=p\Delta V=400\ J$.
		\end{itemize}
	}
\end{vidu}

\loai{Bài toán vận dụng nguyên lí I}
\begin{vidu}%[1P6-3.2K][Loai5] 
	Nén đẳng áp một khối khí ở áp suất 500 kPa làm cho thể tích của nó thay đổi 4 lít. Khối khí truyền ra bên ngoài một nhiệt lượng 1200 J. Độ biến thiên nội năng của khối khí bằng
	\choice
	{1200 J}
	{2000 J}
	{\True 800 J}
	{3200 J}
	\loigiai{
		\begin{itemize}
			\item Khí truyền nhiệt: $Q=-1200\ J$.
			\item Nén khí, thể tích giảm: $A=-p.\Delta V=-5.10^5.(-4.10^{-3})=2000\ J$.
			\item Độ biến thiên nội năng của khí: $\Delta U=A+Q=800\ J$.
	\end{itemize}}
\end{vidu}

\loai{Bài toán giọt thuỷ ngân trong ống}
\begin{vidu} %[1P5-4.2K][Loai4]
	immini[thm]{Ống thủy tinh tiết diện S một đầu kín (hình vẽ), một đầu ngăn bởi giọt thủy ngân. Chiều dài cột không khí bên trong ống thủy tinh là $\ell _1 = 20\ cm$, nhiệt độ bên trong ống là $27\doC$. Coi quá trình biến đổi trạng thái với áp suất là không đổi. Chiều cao của cột không khí bên trong ống khi nhiệt độ tăng thêm $10\doC$ là
		\motcot
		{27,4 cm}
		{12,6 cm}
		{19,33 cm}
		{\True 20,67 cm}}{\begin{tikzpicture}[scale=1,thick,>=stealth']
			\draw[fill=gray!40,opacity=0.5](0,0)node[below left,opacity=1]{$\ell_0$}--++(270:0.5)--++(0:1)--++(90:0.5)--cycle;
			\draw(0,1)--++(270:4)--++(0:1)--++(90:4);
			\draw[<->] (1.3,-0.5)--node[midway,right]{$\ell_1$}(1.3,-3);
	\end{tikzpicture}}  
	\loigiai{
		\begin{itemize}
			\item Trạng thái 1: $T_1 = 300\ K$, $V_1 =\ \ell _1S$.
			\item Trạng thái 2: $T_2 = 310\ K$, $V_2 =\ \ell _2S$.
			\item Quá trình là đẳng áp: $\dfrac{\ell _1S}{300} = \dfrac{\ell _2S}{310} \Rightarrow\ \ell _2 = 20,67\ cm$.
		\end{itemize}
	}
\end{vidu}
\begin{vidu} %[1P5-4.2K][Loai4]
	immini[thm]{Cho áp kế như hình vẽ. Tiết diện ống là $0,1\ cm^2$, biết ở $0\doC$ giọt thủy ngân cách A 30\ cm, ở $5\doC$ giọt thủy ngân cách A 50 cm. Thể tích của bình là  		
		\choice
		{$106,2\ cm^3$}
		{$130\ cm^3$}
		{$250\ cm^3$}
		{\True $106,5\ cm^3$}}
	{\begin{tikzpicture}[scale=1,thick,>=stealth']
			\draw[rounded corners=0.7mm,fill=blue,opacity=0.5](-1,0)--++(180:0.5)--++(270:0.15)--++(0:0.5)--cycle;
			\draw[rounded corners](0,0)--++(180:3)node[above right]{$A$}arc[x radius=1,y radius=1,start
			angle=5,delta angle=350]--++(0:3);
			\node at (-4,0){$V$};
	\end{tikzpicture}}
	\loigiai{
		\Binhluan{
			\begin{itemize}
				\item Trạng thái 1: $T_1 = 273\ K,\ V_1 = V + 0,1.30 = V + 3$. 
				\item Trạng thái 2: $T_2 = 278\ K,\ V_2 = V + 0,1.50 = V + 5$.
				\item Áp dụng định luật Charles: $\dfrac{V + 3}{273} = \dfrac{V + 5}{278} \Rightarrow V = 106,2\ cm^3$.
		\end{itemize}}
		{Áp kế là dụng cụ để đo áp suất}
	}
\end{vidu}
\loai{Bài toán vận tốc khí}
\begin{vidu}%[1P5-4.2G][Loai5]
	Khí từ lò đốt thoát ra theo ống khói hình trụ, ở đầu dưới khí có nhiệt độ $700\doC$ và chuyển động với vận tốc 5 cm/s. Coi áp suất không đổi.
	Tính vận tốc của khí khi lên đến đầu trên ống khói ở nhiệt độ $200\doC$.
	\choice
	{\True 2,43 cm/s}
	{1,43 cm/s}
	{17,50 cm/s}
	{10,29 cm/s}
	\loigiai{
		\begin{itemize}
			\item Gọi $v_1,\ v_2$ là vận tốc của khí ở đầu dưới và đầu trên của ống khói. S là tiết diện của ống khói.
			\item Xét trong khoảng thời gian $t= 1$ s ta có cùng một lượng khí vào đầu dưới và thoát ra khỏi đầu trên: 
			\begin{align*}
				\begin{cases}
					\ct{ở đầu dưới: } T_1 = 973\ K,\ V_1 = v_1S\\ 
					\ct{ở đầu trên: } T_2 = 473\ K,\ {V_2} = v_2S.
				\end{cases}
			\end{align*}
			\item Áp dụng định luật Gay Luy-xác: 
			\begin{align*}
				\dfrac{V_1}{T_1}=\dfrac{V_2}{T_2}\Rightarrow \dfrac{v_1S}{T_1} = \dfrac{v_2S}{T_2} \Rightarrow \dfrac{5}{973} = \dfrac{v_2}{473} \Rightarrow v_2 = 2,43\ cm/s.
			\end{align*}
		\end{itemize}
	}
\end{vidu}
